% Source: https://artofproblemsolving.com/community/c3h1590039_common_problem_contradiction
% Stable identity: aops-topic-1590039-q01
\textbf{(a)} By symmetry, $A$'s partner is equally likely to be any of the other $19$ people, so the probability that it is $B$ is $\boxed{\dfrac{1}{19}}$.

As a count: a pairing of $2m$ people can be built by pairing the lowest-numbered unpaired person with one of the remaining people, repeatedly, so there are $19\cdot17\cdot15\cdots3\cdot1$ pairings of $20$ people. Those with $\{A,B\}$ as a pair correspond to pairings of the other $18$ people, $17\cdot15\cdots1$ of them. The ratio is $\dfrac{1}{19}$.

\textbf{(b)} The student's fraction $\dfrac{2}{20\cdot19}=\dfrac{1}{190}$ is the probability that \emph{one randomly chosen pair of people} is $\{A,B\}$. But a pairing does not consist of one pair; it consists of $10$ pairs, and the event "$A$ and $B$ are paired" happens if \emph{any} of the $10$ pairs is $\{A,B\}$. If we pick one of the $10$ pairs of a random pairing at random, it is a uniformly random pair of people, so it equals $\{A,B\}$ with probability $\dfrac{1}{190}$. These $10$ events (pair number $1$ is $\{A,B\}$, \dots, pair number $10$ is $\{A,B\}$) cannot happen together, so
\[
P(A\text{ and }B\text{ paired})=10\cdot\dfrac{1}{190}=\dfrac{1}{19},
\]
in agreement with (a). The student's denominator counts possible pairs, but the numerator should count the favourable \emph{pairs among those actually formed}, of which there are $10$, not $1$.
